\chapter{Progettazione Logica}

\section{Introduzione e Metodologia di Progettazione Logica}
La progettazione logica ha l'obiettivo di trasformare lo schema concettuale E/R in uno schema relazionale normalizzato, efficiente e direttamente implementabile sul motore relazionale \textbf{MySQL 8.0 (InnoDB)}.

La metodologia accademica del corso prevede i seguenti passaggi formali:
\begin{enumerate}
    \item Stima dei volumi dei dati a regime.
    \item Definizione della tabella delle operazioni e stima delle frequenze d'esecuzione.
    \item Elaborazione degli schemi di navigazione e delle \textbf{tabelle degli accessi} con calcolo del costo computazionale (assegnando peso $L=1$ alle letture e $S=2$ alle scritture).
    \item \textbf{Analisi quantitativa e critica delle ridondanze}, confrontando i costi degli schemi con e senza dato derivato.
    \item Raffinamento dello schema E/R (eliminazione gerarchie IS-A, rimozione identificatori esterni).
    \item Definizione dello schema relazionale finale e traduzione formale delle interrogazioni SQL.
\end{enumerate}

\section{Stima dei Volumi dei Dati}
Si assume uno scenario operativo reale per un'agenzia consolidata attiva sul mercato da \textbf{5 anni}, con un flusso medio a regime di \textbf{40 matrimoni all'anno} (200 matrimoni storicizzati nel sistema). La Tabella~\ref{tab:volumi} illustra i volumi stimati per le 33 entità del dominio.

\begin{table}[h!]
\centering
\small
\begin{tabularx}{\textwidth}{l p{2.5cm} r X}
\toprule
\textbf{Concetto / Entità} & \textbf{Tipologia} & \textbf{Volume (Tuple)} & \textbf{Criteri di Stima e Motivazione} \\
\midrule
\texttt{PERSONA} & Forte (Radice IS-A) & 26.000 & 500 Coniugi, 10 Dipendenti, 25.490 Invitati unici. \\
\texttt{COPPIA} & Forte & 250 & 200 matrimoni chiusi + 50 coppie in trattativa/pianificazione. \\
\texttt{DIPENDENTE} & Specializzazione & 10 & Personale di staff dell'agenzia. \\
\texttt{RUOLO} & Forte & 4 & Mansioni contrattuali codificate. \\
\texttt{RUOLO\_DIPENDENTE} & Reificazione $N:M$ & 16 & Media di 1.6 ruoli per dipendente. \\
\texttt{MOTIVO\_ASSENZA} & Forte & 8 & Categorie assenza (ferie, malattia, congedo, ecc.). \\
\texttt{DISPONIBILITA} & Debole da Dip. & 400 & Storico assenze (8 periodi/anno per dipendente in 5 anni). \\
\texttt{ASSEGNAZIONE\_DIP.} & Reificazione $N:M$ & 320 & Storico planner assegnati alle coppie (con riassegnazioni). \\
\texttt{EVENTO} & Debole da Coppia & 200 & 40 matrimoni/anno $\times$ 5 anni. \\
\texttt{LOCATION} & Forte & 25 & Strutture convenzionate (ville, castelli, casali). \\
\texttt{SVOLGIMENTO\_EVENTO}& Reificazione $N:M$ & 400 & Media di 2 location utilizzate per evento (rito + ricevimento). \\
\texttt{TAVOLO} & Debole da Evento & 2.400 & Media di 12 tavoli per ricevimento. \\
\texttt{BOMBONIERA} & Forte & 12 & Tipologie bomboniera a catalogo. \\
\texttt{CATEGORIA\_FORNITORE}& Forte & 8 & Settori merceologici (catering, foto, musica, fiori, abiti). \\
\texttt{FORNITORE} & Forte & 60 & Imprese partner accreditate (7--8 partner per categoria). \\
\texttt{COLLABORAZIONE} & Debole da Forn. & 120 & Contratti biennali di convenzione rinnovati. \\
\texttt{PREVENTIVO} & Debole (1:N Evento)& 1.200 & Media di 6 preventivi richiesti per matrimonio. \\
\texttt{CONTRATTO} & Debole (1:1 Prev.) & 600 & Tasso di conversione del 50\% (3 contratti stipulati per evento). \\
\texttt{TIMELINE} & Debole (1:1 Evento)& 200 & Esattamente 1 timeline di regia oraria per evento. \\
\texttt{ATTIVITA\_TIMELINE} & Debole da Timeline & 2.400 & Media di 12 momenti scaletta oraria per evento. \\
\texttt{TEMPLATE\_CHECKLIST}& Forte & 6 & Modelli standard per tipologia di rito. \\
\texttt{VOCE\_TEMPLATE} & Debole da Template & 180 & Media di 30 compiti codificati per template. \\
\texttt{CHECKLIST} & Debole (1:1 Evento)& 200 & Esattamente 1 checklist operativa per evento. \\
\texttt{ATTIVITA\_CHECKLIST}& Debole da Check. & 6.000 & 30 compiti operativi monitorati per matrimonio. \\
\texttt{LISTA\_NOZZE} & Debole (1:1 Evento)& 200 & Esattamente 1 lista nozze per evento. \\
\texttt{REGALO} & Debole da Lista & 1.600 & Media di 8 articoli o quote desiderate inserite per lista. \\
\texttt{INVITATO} & Specializzazione & 25.000 & Anagrafica globale degli ospiti invitati. \\
\texttt{PARTECIPAZIONE} & Reificazione $N:M$ & 25.500 & Media di ~127 invitati a matrimonio (alcuni ricorrenti). \\
\texttt{QUOTA\_REGALO} & Reificazione $N:M$ & 10.000 & Circa 50 donazioni/quote versate per evento. \\
\texttt{COMUNICAZIONE} & Debole da Coppia & 3.000 & Media di 12 comunicazioni formali archiviate per coppia. \\
\texttt{APPUNTAMENTO} & Debole da Coppia & 1.250 & Media di 5 incontri di consulenza per coppia. \\
\texttt{RECENSIONE} & Debole (1:1 Evento)& 180 & Tasso di rilascio recensione del 90\% su eventi conclusi. \\
\texttt{FOTO} & Debole da Evento & 1.800 & Media di 9 fotografie ufficiali a portfolio per matrimonio. \\
\bottomrule
\end{tabularx}
\caption{Tavola dei volumi stimati per le entità del dominio a regime quinquennale.}
\label{tab:volumi}
\end{table}

\section{Tabella delle Operazioni e Frequenze}
Le operazioni del sistema vengono classificate in \textbf{Interattive (I)} per le transazioni OLTP di routine e \textbf{Batch (B)} per le elaborazioni analitiche direzionali (Tabella~\ref{tab:operazioni}).

\begin{table}[h!]
\centering
\small
\begin{tabularx}{\textwidth}{l p{7.5cm} c r}
\toprule
\textbf{Codice} & \textbf{Descrizione Operazione} & \textbf{Tipo} & \textbf{Frequenza ($F_{\text{anno}}$)} \\
\midrule
\texttt{OP-01} & Registrazione nuova coppia con due coniugi e assegnazione wedding planner & I & 24 / anno \\
\texttt{OP-02} & Creazione evento con budget ed associazione location e fasce orarie & I & 24 / anno \\
\texttt{OP-03} & Istanziazione automatica checklist operativa a partire dal template di rito & I & 24 / anno \\
\texttt{OP-04} & Inserimento fornitore nel catalogo e associazione categoria & I & 12 / anno \\
\texttt{OP-05} & Emissione richiesta preventivo a fornitore per un matrimonio & I & 240 / anno \\
\texttt{OP-06} & Accettazione preventivo, stipula contratto e registrazione acconto & I & 120 / anno \\
\texttt{OP-07} & Censimento invitato e registrazione partecipazione (RSVP, allergie, menù) & I & 2.400 / anno \\
\texttt{OP-08} & Conferma presenza invitato ed assegnazione posto al tavolo & I & 1.200 / anno \\
\texttt{OP-09} & Inserimento nuovo articolo regalo nella lista nozze & I & 360 / anno \\
\texttt{OP-10} & Registrazione donazione/quota regalo con messaggio auguri da invitato & I & 1.800 / anno \\
\texttt{OP-11} & Chiusura task checklist con data completamento effettiva & I & 1.200 / anno \\
\texttt{OP-12} & Inserimento voce nella timeline di regia oraria con responsabile & I & 480 / anno \\
\texttt{OP-13} & Registrazione recensione e caricamento foto portfolio a evento concluso & I & 36 / anno \\
\midrule
\texttt{STAT-01}& Fornitori maggiormente contrattualizzati per categoria e spesa media & B & 12 / anno (mensile) \\
\texttt{STAT-02}& Monitoraggio scostamento budget (Variance Analysis) su eventi attivi & B & 48 / anno (settim.) \\
\texttt{STAT-03}& Tasso di conversione preventivi in contratti per fornitore & B & 12 / anno (mensile) \\
\texttt{STAT-04}& Analisi saturazione e capienza tavoli per ricevimento & B & 48 / anno (settim.) \\
\texttt{STAT-05}& Avanzamento checklist operativa e tassi di ritardo per fase temporale & B & 48 / anno (settim.) \\
\texttt{STAT-06}& Distribuzione regali e stato di progressione lista nozze & B & 24 / anno \\
\texttt{STAT-07}& Profilassi allergenica e statistiche menù catering & B & 48 / anno (settim.) \\
\texttt{STAT-08}& Ospiti ricorrenti e generosità donazioni su matrimoni multipli & B & 24 / anno \\
\bottomrule
\end{tabularx}
\caption{Tavola delle operazioni principali e relative frequenze stimate.}
\label{tab:operazioni}
\end{table}

\section{Schemi di Navigazione e Tabelle degli Accessi}
Il calcolo formale del costo computazionale adotta le seguenti convenzioni didattiche:
\begin{itemize}[noitemsep]
    \item Ogni accesso in Lettura ($L$) equivale a un costo unitario di \textbf{1}.
    \item Ogni accesso in Scrittura ($S$), che comporta inserimento, cancellazione o modifica su disco con aggiornamento degli indici e log transazionale, equivale a un costo unitario di \textbf{2}.
    \item Costo Annuo $= \text{Costo Singola Esecuzione} \times F_{\text{anno}}$.
\end{itemize}

\subsection{Operazione OP-08: Conferma RSVP e Assegnazione Tavolo}
\textbf{Descrizione:} Dato un invitato (CF) e un matrimonio, si legge la partecipazione, si verifica la capienza residua del tavolo designato e si aggiorna lo stato su \texttt{'CONFERMATO'} assegnando il tavolo. Frequenza $F = 1.200 \text{ volte/anno}$.

\begin{table}[h!]
\centering
\small
\begin{tabular}{l c c c c l}
\toprule
\textbf{Entità / Tabella} & \textbf{Tipo Accesso} & \textbf{Accessi} & \textbf{Costo Unitario} & \textbf{Costo Parziale} & \textbf{Note} \\
\midrule
\texttt{PARTECIPAZIONE} & L & 1 & 1 & 1 & Lettura riga ospite per l'evento \\
\texttt{TAVOLO} & L & 1 & 1 & 1 & Lettura capienza massima tavolo \\
\texttt{PARTECIPAZIONE} & L & 10 & 1 & 10 & Lettura ospiti già assegnati al tavolo \\
\texttt{PARTECIPAZIONE} & S & 1 & 2 & 2 & Scrittura stato confermato e tavolo \\
\midrule
\textbf{Totale Singolo} & & & & \textbf{14} & \\
\textbf{Totale Annuo} & & & & \textbf{16.800} & $14 \times 1.200 \text{ esecuzioni/anno}$ \\
\bottomrule
\end{tabular}
\caption{Tabella degli accessi per l'operazione \texttt{OP-08}.}
\end{table}

\subsection{Operazione OP-10: Registrazione Quota Regalo da Invitato}
\textbf{Descrizione:} Un invitato effettua un versamento per la lista nozze. Il sistema accerta l'identità dell'ospite, verifica la presenza dell'invito per quel matrimonio e registra la quota monetaria. Frequenza $F = 1.800 \text{ volte/anno}$.

\begin{table}[h!]
\centering
\small
\begin{tabular}{l c c c c l}
\toprule
\textbf{Entità / Tabella} & \textbf{Tipo Accesso} & \textbf{Accessi} & \textbf{Costo Unitario} & \textbf{Costo Parziale} & \textbf{Note} \\
\midrule
\texttt{INVITATO} & L & 1 & 1 & 1 & Controllo anagrafica ospite \\
\texttt{PARTECIPAZIONE} & L & 1 & 1 & 1 & Verifica invito valido per l'evento \\
\texttt{REGALO} & L & 1 & 1 & 1 & Lettura prezzo target (se regalo specifico) \\
\texttt{QUOTA\_REGALO} & L & 6 & 1 & 6 & Lettura quote già versate per l'articolo \\
\texttt{QUOTA\_REGALO} & S & 1 & 2 & 2 & Inserimento record quota regalo \\
\midrule
\textbf{Totale Singolo} & & & & \textbf{11} & \\
\textbf{Totale Annuo} & & & & \textbf{19.800} & $11 \times 1.800 \text{ esecuzioni/anno}$ \\
\bottomrule
\end{tabular}
\caption{Tabella degli accessi per l'operazione \texttt{OP-10}.}
\end{table}

\subsection{Operazione STAT-02: Monitoraggio Scostamento Budget (Variance Analysis)}
\textbf{Descrizione:} Elaborazione analitica settimanale ($F = 48 \text{ volte/anno}$) che analizza i 40 matrimoni attivi in pianificazione, sommando contratti fornitore e costi di affitto location per confrontarli con il budget stanziato.

\begin{table}[h!]
\centering
\small
\begin{tabular}{l c c c c l}
\toprule
\textbf{Entità / Tabella} & \textbf{Tipo Accesso} & \textbf{Accessi} & \textbf{Costo Unitario} & \textbf{Costo Parziale} & \textbf{Note} \\
\midrule
\texttt{EVENTO} & L & 40 & 1 & 40 & Scansione dei 40 eventi attivi \\
\texttt{PREVENTIVO} & L & 240 & 1 & 240 & 6 preventivi medi per evento ($40 \times 6$) \\
\texttt{CONTRATTO} & L & 120 & 1 & 120 & 3 contratti stipulati per evento ($40 \times 3$) \\
\texttt{SVOLGIMENTO\_EVENTO} & L & 80 & 1 & 80 & 2 location per evento ($40 \times 2$) \\
\texttt{LOCATION} & L & 80 & 1 & 80 & Lettura costo affitto location \\
\midrule
\textbf{Totale Singolo} & & & & \textbf{560} & \\
\textbf{Totale Annuo} & & & & \textbf{26.880} & $560 \times 48 \text{ esecuzioni/anno}$ \\
\bottomrule
\end{tabular}
\caption{Tabella degli accessi per l'operazione \texttt{STAT-02}.}
\end{table}

\section{Analisi Critica delle Ridondanze}
In conformità alle direttive del corso, sono stati sottoposti a rigorosa analisi quantitativa quattro attributi candidati di ridondanza, confrontando sistematicamente lo Schema $S_1$ (senza ridondanza) e lo Schema $S_2$ (con ridondanza).

\subsection{Ridondanza 1 (ACCETTATA): Totale Raccolto in LISTA\_NOZZE}
\begin{itemize}[noitemsep]
    \item \textbf{Dato Derivabile:} $\text{LISTA\_NOZZE.totale\_raccolto} = \sum_{q \in \text{QUOTA\_REGALO}} q.\text{importo}$.
    \item \textbf{Operazioni coinvolte:} 
    \begin{itemize}[noitemsep]
        \item $OP_A$: Registrazione quota regalo (\texttt{OP-10}, $F_A = 1.800 \text{ volte/anno}$).
        \item $OP_B$: Visualizzazione lista nozze da portale sposi o invitati ($F_B = 5.400 \text{ volte/anno}$).
    \end{itemize}
\end{itemize}

\textbf{Calcolo dei Costi Comparativi:}
\begin{itemize}
    \item \textbf{Schema $S_1$ (Senza Ridondanza):}
    \begin{itemize}[noitemsep]
        \item Scrittura ($OP_A$): Inserimento in \texttt{QUOTA\_REGALO} ($1 \text{ scrittura } = 2$). Costo annuo: $2 \times 1.800 = \mathbf{3.600}$.
        \item Lettura ($OP_B$): Scansione delle 50 quote versate in media per evento ($50 \text{ letture } = 50$). Costo annuo: $50 \times 5.400 = \mathbf{270.000}$.
        \item \textbf{Costo Totale Annuo $S_1$:} $3.600 + 270.000 = \mathbf{273.600}$.
    \end{itemize}
    \item \textbf{Schema $S_2$ (Con Ridondanza):}
    \begin{itemize}[noitemsep]
        \item Scrittura ($OP_A$): Inserimento quota ($2$) + lettura tupla lista ($1$) + aggiornamento accumulatore lista ($2$) = $5$. Costo annuo: $5 \times 1.800 = \mathbf{9.000}$.
        \item Lettura ($OP_B$): Accesso diretto al valore precalcolato nella tupla \texttt{LISTA\_NOZZE} ($1 \text{ lettura } = 1$). Costo annuo: $1 \times 5.400 = \mathbf{5.400}$.
        \item \textbf{Costo Totale Annuo $S_2$:} $9.000 + 5.400 = \mathbf{14.400}$.
    \end{itemize}
\end{itemize}

\[ \Delta \text{Costo} = 273.600 - 14.400 = \mathbf{+259.200 \text{ accessi/anno risparmiati (-94.7\%)}} \]

\textbf{Decisione Formale: ACCETTATA.} Il risparmio computazionale è imponente; la sincronizzazione è garantita tramite Trigger atomico \texttt{trg\_quota\_aggiorna\_totale}.

\subsection{Ridondanza 2 (SCARTATA): Costo Totale Impegnato in EVENTO}
\begin{itemize}[noitemsep]
    \item \textbf{Dato Derivabile:} Somma dei contratti stipulati e canoni location:
    $\text{CostoImpegnato} = \sum c.\text{importo\_concordato} + \sum l.\text{costo\_affitto}$.
    \item \textbf{Operazioni coinvolte:} Modifica contrattuale/location ($F = 180 \text{ volte/anno}$) e reportistica budget ($F = 148 \text{ volte/anno}$).
\end{itemize}

\textbf{Valutazione Comparativa:}
Sebbene l'aritmetica evidenzi un risparmio teorico ($29.960$ vs $6.820$), \textbf{la ridondanza viene respinta per ragioni di architettura}:
\begin{enumerate}[noitemsep]
    \item L'attributo dipende da due tabelle distinte (\texttt{CONTRATTO} e \texttt{SVOLGIMENTO\_EVENTO}), moltiplicando i trigger necessari e aumentando il rischio di lock concorrenti durante le stipule.
    \item La lettura avviene quasi esclusivamente in una procedura batch settimanale, dove una query aggregata garantisce assoluta consistenza finanziaria senza alcun impatto percettibile.
\end{enumerate}

\subsection{Ridondanza 3 (SCARTATA): Posti Occupati in TAVOLO}
\begin{itemize}[noitemsep]
    \item \textbf{Dato Derivabile:} $\text{COUNT}(p \in \text{PARTECIPAZIONE} : p.\text{tavolo} = t)$.
    \item \textbf{Valutazione:} Mantenere un contatore denormalizzato su \texttt{TAVOLO} comporterebbe un lock esclusivo sulla riga del tavolo ad ogni assegnazione posto. Poiché ciascun tavolo ospita mediamente solo 10 persone, la query \texttt{COUNT(*)} indicizzata impegna meno di 0.1 ms. La ridondanza viene **SCARTATA** per non creare colli di bottiglia di concorrenza nella GUI.
\end{itemize}

\subsection{Ridondanza 4 (ACCETTATA): Flag Completato in REGALO}
\begin{itemize}[noitemsep]
    \item \textbf{Dato Derivabile:} Flag \texttt{'S'/'N'} indicante se la somma delle quote copre il prezzo: $\sum q.\text{importo} \ge r.\text{prezzo}$.
    \item \textbf{Valutazione:} Durante la navigazione della lista regali da parte degli invitati ($5.400 \text{ accessi/anno}$), filtrare i regali completati senza il flag richiederebbe una complessa \texttt{LEFT JOIN} con \texttt{GROUP BY} e \texttt{HAVING} su tutte le quote (costo annuo $316.800$). Con il flag ridondante indicizzato, la lettura costa $43.200$. Risparmio: $\mathbf{+269.388 \text{ accessi/anno (-85.0\%)}}$.
    \item \textbf{Decisione Formale: ACCETTATA}, aggiornata dal trigger delle quote.
\end{itemize}

\begin{table}[h!]
\centering
\small
\begin{tabularx}{\textwidth}{l p{3.8cm} c c X}
\toprule
\textbf{Ridondanza} & \textbf{Entità Destinazione} & \textbf{Costo $S_1$} & \textbf{Costo $S_2$} & \textbf{Esito e Motivazione} \\
\midrule
\texttt{totale\_raccolto} & \texttt{LISTA\_NOZZE} & 273.600 & 14.400 & \textbf{ACCETTATA:} Riduce gli accessi del 94.7\% in visualizzazione lista. \\
\addlinespace
\texttt{costo\_impegnato} & \texttt{EVENTO} & 29.960 & 6.820 & \textbf{SCARTATA:} Evita trigger multipli e rischi di disallineamento contabile. \\
\addlinespace
\texttt{posti\_occupati} & \texttt{TAVOLO} & 14.400 & 10.800 & \textbf{SCARTATA:} Previene lock di riga concorrenti durante il seating planning. \\
\addlinespace
\texttt{completato} & \texttt{REGALO} & 316.800 & 47.412 & \textbf{ACCETTATA:} Elimina pesanti aggregazioni runtime nel catalogo regali. \\
\bottomrule
\end{tabularx}
\caption{Quadro sinottico dell'analisi quantitativa delle ridondanze.}
\label{tab:quadro_ridondanze}
\end{table}

\section{Raffinamento e Ristrutturazione dello Schema E/R}
\subsection{Eliminazione delle Gerarchie di Generalizzazione (IS-A)}
La gerarchia avente come radice \texttt{PERSONA} e figli \texttt{DIPENDENTE} e \texttt{INVITATO} è stata ristrutturata come segue:
\begin{itemize}
    \item \textbf{Sottotipo INVITATO:} Possiede pochi attributi specifici (\texttt{rapporto\_parentela}, \texttt{nucleo\_familiare}) e partecipa ad associazioni ad altissima frequenza. Si è adottato il \textbf{collasso verso l'alto parziale / sostituzione con associazione}: \texttt{PERSONA} conserva le generalità anagrafiche, mentre la tabella \texttt{INVITATO} estende \texttt{PERSONA} tramite FK primaria su \texttt{codice\_fiscale}.
    \item \textbf{Sottotipo DIPENDENTE:} Ha numerose associazioni esclusive (\texttt{RUOLO\_DIPENDENTE}, \texttt{DISPONIBILITA}, \texttt{ASSEGNAZIONE\_DIPENDENTE}). È stata adottata la \textbf{sostituzione con associazione}: entità autonoma \texttt{DIPENDENTE} con PK \texttt{codice\_fiscale} in integrità referenziale su \texttt{PERSONA}.
\end{itemize}

\subsection{Eliminazione degli Identificatori Esterni}
Le entità deboli concettuali sono state tradotte importando le chiavi primarie delle entità proprietarie:
\begin{itemize}[noitemsep]
    \item \texttt{EVENTO:} Assorbe la chiave della coppia $\rightarrow \text{PK} = \{\underline{\text{username\_coppia}, \text{data}}\}$.
    \item \texttt{TAVOLO:} Debole da evento $\rightarrow \text{PK} = \{\underline{\text{username\_coppia}, \text{data\_evento}, \text{nome\_numero}}\}$.
    \item \texttt{LISTA\_NOZZE, CHECKLIST, TIMELINE, RECENSIONE:} Relazioni 1:1 con evento $\rightarrow \text{PK} = \{\underline{\text{username\_coppia}, \text{data\_evento}}\}$.
    \item \texttt{PREVENTIVO:} Debole da fornitore $\rightarrow \text{PK} = \{\underline{\text{fornitore\_piva}, \text{numero\_preventivo}}\}$ con FK verso evento.
\end{itemize}

\section{Schema Relazionale Finale}
Si formalizza lo schema logico relazionale (chiavi primarie sottolineate, chiavi esterne con asterisco):

\small
\begin{enumerate}[noitemsep]
    \item $\text{PERSONA}(\underline{\text{codice\_fiscale}}, \text{nome}, \text{cognome}, \text{email}, \text{telefono})$
    \item $\text{COPPIA}(\underline{\text{username}}, \text{data\_prevista}, \text{budget\_indicativo}, \text{cf\_coniuge1}^*, \text{cf\_coniuge2}^*)$ \\
    \textit{FK: $\text{cf\_coniuge1}, \text{cf\_coniuge2} \rightarrow \text{PERSONA}$; CHECK: $\text{cf\_coniuge1} < \text{cf\_coniuge2}$}
    \item $\text{DIPENDENTE}(\underline{\text{codice\_fiscale}}^*, \text{email\_aziendale}, \text{data\_assunzione}, \text{data\_fine\_rapporto}, \text{stato})$ \\
    \textit{FK: $\text{codice\_fiscale} \rightarrow \text{PERSONA}$}
    \item $\text{RUOLO}(\underline{\text{nome\_ruolo}}, \text{descrizione}, \text{max\_eventi\_contemporanei})$
    \item $\text{RUOLO\_DIPENDENTE}(\underline{\text{dipendente\_cf}}^*, \underline{\text{nome\_ruolo}}^*, \text{data\_assegnazione}, \text{principale})$
    \item $\text{MOTIVO\_ASSENZA}(\underline{\text{categoria}}, \text{descrizione}, \text{remunerata})$
    \item $\text{DISPONIBILITA}(\underline{\text{dipendente\_cf}}^*, \underline{\text{data\_inizio}}, \text{data\_fine}, \text{tipo}, \text{stato}, \text{categoria\_assenza}^*)$
    \item $\text{ASSEGNAZIONE\_DIPENDENTE}(\underline{\text{username\_coppia}}^*, \underline{\text{dipendente\_cf}}^*, \underline{\text{data\_inizio}}, \text{data\_fine}, \text{motivo\_cambio}, \text{attiva})$
    \item $\text{EVENTO}(\underline{\text{username\_coppia}}^*, \underline{\text{data}}, \text{budget\_effettivo}, \text{n\_invitati\_previsti}, \text{tipo\_cerimonia}, \text{tema\_stile}, \text{stato})$
    \item $\text{LOCATION}(\underline{\text{nome\_location}}, \text{indirizzo}, \text{citta}, \text{capienza\_max}, \text{costo\_affitto}, \text{referente}, \text{telefono})$
    \item $\text{SVOLGIMENTO\_EVENTO}(\underline{\text{username\_coppia}}^*, \underline{\text{data\_evento}}^*, \underline{\text{nome\_location}}^*, \underline{\text{fascia\_oraria}}, \text{tipo\_utilizzo})$
    \item $\text{TAVOLO}(\underline{\text{username\_coppia}}^*, \underline{\text{data\_evento}}^*, \underline{\text{nome\_numero}}, \text{capienza}, \text{forma}, \text{zona\_sala})$
    \item $\text{BOMBONIERA}(\underline{\text{tipologia}}, \text{descrizione}, \text{costo\_unitario})$
    \item $\text{CATEGORIA\_FORNITORE}(\underline{\text{nome\_categoria}}, \text{descrizione})$
    \item $\text{FORNITORE}(\underline{\text{partita\_iva}}, \text{ragione\_sociale}, \text{indirizzo}, \text{telefono}, \text{email}, \text{attivo}, \text{nome\_categoria}^*)$
    \item $\text{COLLABORAZIONE}(\underline{\text{fornitore\_piva}}^*, \underline{\text{data\_inizio}}, \text{data\_fine}, \text{stato}, \text{commissione\_percentuale})$
    \item $\text{PREVENTIVO}(\underline{\text{fornitore\_piva}}^*, \underline{\text{numero\_preventivo}}, \text{data\_emissione}, \text{data\_scadenza}, \text{importo}, \text{stato}, \text{username\_coppia}^*, \text{data\_evento}^*)$
    \item $\text{CONTRATTO}(\underline{\text{fornitore\_piva}}^*, \underline{\text{numero\_preventivo}}^*, \text{data\_firma}, \text{importo\_concordato}, \text{acconto\_versato}, \text{stato\_pagamento})$
    \item $\text{TIMELINE}(\underline{\text{username\_coppia}}^*, \underline{\text{data\_evento}}^*, \text{note\_generali})$
    \item $\text{ATTIVITA\_TIMELINE}(\underline{\text{username\_coppia}}^*, \underline{\text{data\_evento}}^*, \underline{\text{orario\_previsto}}, \text{titolo\_attivita}, \text{ora\_inizio}, \text{ora\_fine}, \text{dipendente\_cf}^*)$
    \item $\text{TEMPLATE\_CHECKLIST}(\underline{\text{id\_template}}, \text{nome\_template}, \text{descrizione}, \text{tipo\_rito})$
    \item $\text{VOCE\_TEMPLATE}(\underline{\text{id\_template}}^*, \underline{\text{titolo\_voce}}, \text{descrizione}, \text{giorni\_prima\_evento}, \text{priorita}, \text{fase\_temporale})$
    \item $\text{CHECKLIST}(\underline{\text{username\_coppia}}^*, \underline{\text{data\_evento}}^*, \text{id\_template}^*)$
    \item $\text{ATTIVITA\_CHECKLIST}(\underline{\text{username\_coppia}}^*, \underline{\text{data\_evento}}^*, \underline{\text{titolo\_attivita}}, \text{scadenza}, \text{data\_completamento}, \text{priorita}, \text{fase\_temporale})$
    \item $\text{LISTA\_NOZZE}(\underline{\text{username\_coppia}}^*, \underline{\text{data\_evento}}^*, \text{iban}, \text{banca}, \text{intestatario}, \text{totale\_raccolto}^{\text{RID}})$
    \item $\text{REGALO}(\underline{\text{username\_coppia}}^*, \underline{\text{data\_evento}}^*, \underline{\text{descrizione}}, \text{prezzo}, \text{tipo\_frazionamento}, \text{completato}^{\text{RID}})$
    \item $\text{INVITATO}(\underline{\text{codice\_fiscale}}^*, \text{rapporto\_parentela}, \text{nucleo\_familiare}, \text{tipologia\_bomboniera}^*)$
    \item $\text{PARTECIPAZIONE}(\underline{\text{codice\_fiscale}}^*, \underline{\text{username\_coppia}}^*, \underline{\text{data\_evento}}^*, \text{stato}, \text{esigenze\_alimentari}, \text{intolleranze}, \text{menu\_scelto}, \text{tavolo\_nome\_numero}^*)$
    \item $\text{QUOTA\_REGALO}(\underline{\text{codice\_fiscale}}^*, \underline{\text{username\_coppia}}^*, \underline{\text{data\_evento}}^*, \underline{\text{data\_ora}}, \text{importo}, \text{metodo\_pagamento}, \text{stato}, \text{descrizione\_regalo}^*)$
    \item $\text{COMUNICAZIONE}(\underline{\text{id\_comunicazione}}, \text{username\_coppia}^*, \text{data\_ora}, \text{mittente}, \text{canale}, \text{oggetto}, \text{testo})$
    \item $\text{APPUNTAMENTO}(\underline{\text{id\_appuntamento}}, \text{username\_coppia}^*, \text{dipendente\_cf}^*, \text{data\_ora}, \text{luogo}, \text{stato})$
    \item $\text{RECENSIONE}(\underline{\text{username\_coppia}}^*, \underline{\text{data\_evento}}^*, \text{data}, \text{voto}, \text{testo}, \text{tag\_stile})$
    \item $\text{FOTO}(\underline{\text{username\_coppia}}^*, \underline{\text{data\_evento}}^*, \underline{\text{percorso\_file}}, \text{didascalia}, \text{autore}, \text{data\_scatto})$
\end{enumerate}
\normalsize

\section{Interrogazioni SQL Analitiche (STAT-01 .. STAT-08)}
Di seguito si riportano due delle query più significative tra le 8 implementate ed eseguite dal sistema:

\subsection{STAT-01: Fornitori Maggiormente Contrattualizzati e Spesa Media}
\begin{lstlisting}
SELECT 
    f.nome_categoria AS Categoria,
    f.partita_iva AS PartitaIVA,
    f.ragione_sociale AS Fornitore,
    COUNT(c.numero_preventivo) AS NumeroContrattiStipulati,
    COALESCE(SUM(c.importo_concordato), 0.00) AS TotaleFatturato,
    ROUND(COALESCE(AVG(c.importo_concordato), 0.00), 2) AS SpesaMediaPerContratto,
    DENSE_RANK() OVER (
        PARTITION BY f.nome_categoria 
        ORDER BY COUNT(c.numero_preventivo) DESC, SUM(c.importo_concordato) DESC
    ) AS PosizioneClassifica
FROM FORNITORE f
LEFT JOIN PREVENTIVO p 
    ON f.partita_iva = p.fornitore_piva
LEFT JOIN CONTRATTO c 
    ON p.numero_preventivo = c.numero_preventivo 
   AND p.fornitore_piva = c.fornitore_piva
WHERE f.attivo = 'S'
GROUP BY f.nome_categoria, f.partita_iva, f.ragione_sociale
ORDER BY f.nome_categoria, PosizioneClassifica;
\end{lstlisting}

\subsection{STAT-02: Monitoraggio Scostamento Budget (Variance Analysis)}
\begin{lstlisting}
SELECT 
    e.username_coppia AS Coppia,
    e.data AS DataEvento,
    e.budget_effettivo AS BudgetAllocato,
    COALESCE(SUM(DISTINCT l.costo_affitto), 0.00) AS CostoLocation,
    COALESCE(SUM(c.importo_concordato), 0.00) AS CostoFornitori,
    (COALESCE(SUM(DISTINCT l.costo_affitto), 0.00) + COALESCE(SUM(c.importo_concordato), 0.00)) AS TotaleImpegnato,
    (e.budget_effettivo - (COALESCE(SUM(DISTINCT l.costo_affitto), 0.00) + COALESCE(SUM(c.importo_concordato), 0.00))) AS ResiduoMargine,
    ROUND(((COALESCE(SUM(DISTINCT l.costo_affitto), 0.00) + COALESCE(SUM(c.importo_concordato), 0.00)) / e.budget_effettivo) * 100, 2) AS PercentualeConsumata
FROM EVENTO e
LEFT JOIN SVOLGIMENTO_EVENTO se 
    ON e.username_coppia = se.username_coppia AND e.data = se.data_evento
LEFT JOIN LOCATION l 
    ON se.nome_location = l.nome_location
LEFT JOIN PREVENTIVO p 
    ON e.username_coppia = p.username_coppia AND e.data = p.data_evento
LEFT JOIN CONTRATTO c 
    ON p.numero_preventivo = c.numero_preventivo AND p.fornitore_piva = c.fornitore_piva
WHERE e.stato IN ('in_pianificazione', 'confermato', 'completato')
GROUP BY e.username_coppia, e.data, e.stato, e.budget_effettivo
ORDER BY e.data ASC;
\end{lstlisting}
